\section{Java survival components: event clocks and defensible intervention services}
\label{sec:exec24-longitudinal-java}
Java survival components answer different parts of a banking problem. A
collections manager needs to compare action sequences for an eligible borrower
population; a capital analyst may first need event-time prediction. Combining
records held by several institutions introduces a confidentiality problem.
These requirements should not collapse into whether a package implements Cox
regression. Selected primary sources establish four candidates: H2O-3's
distributed fitting and portable scoring, BioJava's survival module,
FRESCO-stat's secure computation, and the distinct repository
\texttt{robbyjo/JLinAlg}. They do not establish exhaustive Java coverage,
qualified integration or a longitudinal causal estimator simply because a
survival routine exists.
\citep{exec24_longjava_h2o,exec24_longjava_biojava,exec24_longjava_fresco,exec24_longjava_jlinalg}

\subsection{Specify the event, clock and target first}
Let $L_i$ be entry time, $T_i$ the end of an observed interval, $D_i$ its
event indicator, and $X_i(t)$ information available before a decision at
$t$. In stratum $s$, stipulate the working model
$\lambda_i(t)=\lambda_{0s}(t)\exp\{X_i(t)^\top\beta+o_i(t)\}$, where
$o_i$ is a declared offset. Its risk set uses $L_i<t\leq T_i$. Neither
a fitted coefficient nor concordance identifies the consequence of changing
collections policy. A causal use still needs the treatment rule, population,
observation mechanism, sequential support and identification assumptions
developed earlier. A confounder affected by earlier treatment does not become
a valid ordinary adjustment variable merely because a start--stop row can
be supplied. The comparison separates the event-time computational component
from the intervention argument. The clock convention is explicit in the
selected H2O and JLinAlg source. \citep{exec24_longjava_h2o,exec24_longjava_jlinalg}

For a fixed baseline profile, write $r=\exp(X^\top\beta+o)$.
Integrating the stipulated hazard gives $S(t\mid X)=S_0(t)^r$.
A log score, relative score and probability have different units. With
$r=2$, baseline survivals $1/2$ and $3/4$ imply risks $3/4$ and $7/16$.
Under a stipulated USD~1,000 loss and USD~600 premium, margins are
$-150$ and $162.50$ dollars. The same relative score permits opposite
decisions. These are exact authored laws, not bank observations or fitted
package outputs. A CLI or MCP response must bind horizon, baseline, stratum
and centering before returning monetary consequences.
\citep{exec24_longjava_reference}

Event coding matters too. In a designed ten-borrower cohort, one default and
four prepayments give observed default fraction $1/10$; collapsing the events
gives $1/2$. Applying the same stipulated USD~1,000 default loss yields
USD~100 versus USD~500 per borrower. This coding witness is not a
counterfactual estimator. Prepayment as a competing event, censoring under an
observation model, and hypothetical elimination of prepayment remain distinct
targets. The earlier R and Python comparisons provide alternatives for
repeated-treatment requests. \citep{exec24_longjava_reference}

\subsection{H2O-3: separate fitting, baseline and scoring contracts}
The selected H2O commit is dated 20 September 2026; its build property says
3.47.0 and \texttt{bleeding\_edge}. This is a pinned development profile,
not a supported-release claim. Its Cox implementation accepts start and stop
columns, strata, offsets and observation weights, with Efron or Breslow ties.
The introductory document calls time columns integer-valued, but inspected
validation requires numeric columns and maps distinct stop times to indices.
The per-row discretizer rejects $L_i\geq T_i$; the preliminary global check
alone is not the full validation path. The common task skips bad or
zero-weight rows, while the processing routine rejects nonpositive weights
that reach it. A wrapper must report exclusions and resolved weights.
\citep{exec24_longjava_h2o}

The event routine subtracts the minimum observed event value and treats a
positive difference as an event. Multilevel codes are therefore not an
automatic competing-risk model. Require an explicit binary event mapping.
Coefficient covariance in the inspected path is the negative inverse Hessian;
the R summary bridge explicitly says robust variance was not used. Supplying
inverse-probability weights alone does not supply uncertainty for estimated
weights, repeated histories or branch shocks. \citep{exec24_longjava_h2o}

The documentation has a weighted-likelihood discrepancy. Its displayed
no-tie contribution is $\exp(w_i\eta_i)/\sum_{j\in R_i}w_j\exp(\eta_j)$,
whereas its Breslow expression and inspected Breslow code raise the denominator
to the event weight. With event weight two, another weight one, and relative
risks one and two, these expressions are $1/4$ and $1/16$. At covariates zero
and one with coefficient $\log 2$, the corresponding log-likelihood
derivatives are $-1/2$ and $-1$. This is not an ignorable constant.
Specify the implementation and tie/weight semantics, retain the discrepancy,
and require an independent fixture before native use.
\citep{exec24_longjava_h2o,exec24_longjava_reference}

Baseline outputs require another check. One cumulative-hazard array uses
tie-specific risk-set denominators. Named baseline frames use a separate loop
starting with total risk and subtracting rows at stop times; that loop does
not apply the Efron correction or add delayed entrants at entry times.
An abstract witness fixes unit risks: A fails at time one, B enters at one
and stops at two, and C is followed from zero to two. The declared risk set
at time one is A and C, giving increment $1/2$; an all-row initial total gives
$1/3$. Two tied events among three unit-risk rows separately give Breslow
$2/3$ and Efron $1/3+1/2=5/6$. These are source-path implications at
supplied risks, not observed H2O runtime failures. Test the exact exported
baseline rather than treating similarly named outputs as interchangeable.
\citep{exec24_longjava_h2o,exec24_longjava_reference}

The scoring frame is named \texttt{lp} and returns a centered linear
predictor. The MOJO scorer likewise returns a linear combination with
stratum centering and offset, not an absolute survival probability. Its
reader loads coefficients, means, strata and interactions; this score alone
does not provide a horizon-specific baseline. The source includes parity
tests for optional columns. One categorical/stratified test deliberately uses
Breslow because of a documented Java~11-plus Efron issue. That comment
identifies a qualification question, not failure in every current environment.
An original H2O fitting service and separately qualified portable scorer are
plausible choices; a whole-library port has no demonstrated benefit here.
\citep{exec24_longjava_h2o}

\subsection{BioJava: an existing port still needs semantic review}
The survival module belongs to a 7.3.1-SNAPSHOT parent. Source describes a
port of portions of R survival and explicitly says not all variants were
implemented. The selected helper accepts one event/censoring time, weights,
strata and clusters. Its row object stores one time; this path does not
establish a general counting-process interface. Controls expose ties,
scaling, tolerances and a nonconvergence flag. The helper skips nonpositive
weights before the downstream weight-use choice, catches exceptions, prints
a trace and returns null. A worker must preserve retained-row identities
and convert null into a typed failure, never a zero coefficient.
\citep{exec24_longjava_biojava}

The robust path saves naive variance, obtains score or dfbeta residuals,
groups by cluster when supplied and forms a cross-product. Strata conversion
still needs care: the fitting arrays describe terminal-row indicators,
whereas the helper reads integers and the comparator sorts by time and status.
Ordinary group labels require a verified mapping. Two branches also need
repair or refusal before reliance. A conjunction in the cluster guard can
dereference a null value before the intended error, while an empty string
passes. In robust unclustered cleanup, the loop saving predictors and scores
does not advance its index; restoration does. For abstract predictors $(2,3)$
and scores $(4,8)$, this structure restores $(3,0)$ and $(8,0)$. The cluster
branch does increment its index. This local trace does not claim all fitted
coefficients or robust variances are wrong: it identifies mutated row state
requiring a native fixture and reviewed repair.
\citep{exec24_longjava_biojava,exec24_longjava_reference}

Cluster identity is part of the sampling contract. For four scalar score
contributions $(1,1,-1,-1)$ and stipulated inverse-information factor $1/4$,
row independence gives $(1/4)^2(1+1+1+1)=1/4$. Grouping the first and last
pairs gives $(1/4)^2(2^2+(-2)^2)=1/2$. This is exact sandwich arithmetic,
not a coverage claim. Borrower clusters do not automatically cover branch
or calendar dependence. \citep{exec24_longjava_reference}

\subsection{FRESCO-stat: confidentiality is a separate contract}
At the pinned March~2023 commit, FRESCO-stat declares 0.6-SNAPSHOT,
FRESCO~1.3.3 and Java~9 source/target settings. Its survival row holds
secret-shared covariates, one time and a field named \texttt{censored}.
The gradient multiplies contributions by that field as an event-status bit;
the inspected example uses one for an observed event. Resolve the name
against behavior rather than guessing its encoding.
\citep{exec24_longjava_fresco}

The public class comment says gradient descent, but the optimizer computes
a Hessian pseudoinverse, multiplies it by the gradient, and subtracts a
scaled step for a fixed iteration count. The result has coefficients,
standard errors and a partial-likelihood-ratio statistic; the inspected loop
has no further convergence criterion. Fixed-point exponentials, inverses
and square roots add numerical obligations. The README warns that overflow
can produce very large outputs. A returned number does not certify
convergence or identified causal uncertainty. \citep{exec24_longjava_fresco}

The sorter identifies tied times using a helper that encrypts values under
a fresh shared key and opens ciphertexts. Equality groups then become public
values used in the risk-set calculation. Lists $(1,1,2)$ and $(10,10,20)$
have the same equality partition: their first two entries match. Revealing
that partition is different from reconstructing raw times, but the disclosure
contract must name it. This inspection proves no complete security theorem
or differential privacy for fitted outputs. Protocol suite, adversary,
modulus, precision, recipients and repeated-query policy remain required.
\citep{exec24_longjava_fresco,exec24_longjava_reference}

A separately qualified secure worker is a candidate when cross-custodian
analysis justifies its overhead. Centralized computation is a comparator
under a different access agreement. Secure computation does not create
treatment support, remove confounding or establish input truth. No upstream
protocol or test was executed in this comparison.
\citep{exec24_longjava_fresco}

\subsection{JLinAlg: explicit diagnostics and a narrower CLI}
The selected \texttt{robbyjo/JLinAlg} commit declares 0.3.6. Its vignette
distinguishes a general right-censored Cox/kinship-frailty CLI from Java APIs
for start--stop rows, strata and other models. The inspected fixed-effect API
validates finite $0\leq L_i<T_i$, nonnegative strata and at least one event.
It prepares risk-set structures and owns a backend in a closeable single-worker
fitter. The optimizer rescales predictors, uses step halving, tests coefficient
and score convergence, and returns convergence status and backend provenance
with covariance and baseline hazards. A result object need not have a true
convergence flag. \citep{exec24_longjava_jlinalg}

The counting-process baseline adds rows whose start is strictly before the
event and removes rows whose stop is before it; Efron and Breslow increments
are explicit. Separate diagnostics form cluster-summed score outer products
between covariance factors. They do not replace the ordinary errors already
in the fit result. The wrapper must name its selected covariance and cluster
population. The inspected fit signature has no case-weight argument; these
sources do not establish a general weighted marginal structural model.
Intervals and robust errors are components, not a sequential intervention
design. \citep{exec24_longjava_jlinalg}

The vignette reports development-host timings and broader frailty tests.
They are not independently reproduced here, and selected fixed-effect review
does not qualify mixed, sparse or GPU paths. The build pins a JDistlib~0.10.2
jar with an expected SHA-256, a useful artifact declaration that still does
not establish a complete locked environment or transitive rights. Backend
selection, helper numerics, runtime behavior and full costs remain open.
\citep{exec24_longjava_jlinalg}

\subsection{Adoption choices and executable next steps}
Direct texts declare Apache~2.0 for H2O, LGPL~2.1 in BioJava's root,
MIT for FRESCO-stat, and GPL~2 in JLinAlg's root with GPL-2.0-or-later
headers in the selected Java files. BioJava's POM says LGPL v2 while linking
the 2.1 text. Preserve these declarations and resolve distribution and
transitive artifacts with a rights owner; process separation is not legal
clearance. \citep{exec24_longjava_h2o,exec24_longjava_biojava,exec24_longjava_fresco,exec24_longjava_jlinalg}

The proposal choice is reference use of all four bounded profiles, followed
conditionally by original-language workers. H2O is a distributed-fitting and
portable-scoring candidate; BioJava offers inspectable robust machinery with
the stated branch and strata obligations; FRESCO-stat addresses a specialized
custodian problem; JLinAlg offers explicit interval and diagnostic APIs.
R/Python workers remain alternatives for repeated-treatment targets.
Reject automatic causal routing from a survival label, automatic log-score
conversion into money, and a whole-library port justified only by language
preference. \citep{exec24_longjava_h2o,exec24_longjava_biojava,exec24_longjava_fresco,exec24_longjava_jlinalg}

Execution starts with the method owner specifying target, clock, event
dictionary, observation model, treatment rule and dependence unit. The
adapter owner freezes a worker artifact and typed input/output schema,
including retained/rejected IDs, resolved options, convergence, covariance,
baseline, warnings and artifact identity. Agents may propose declarations;
unresolved semantics require a structured refusal before fitting. CLI jobs
and MCP calls should preserve the workflow chapter's request, cancellation
and artifact rules rather than silently converting failure into success.

The evaluation owner can preregister generated cases for delayed entry,
ties, nonunit weights, event coding, missing clusters, nonconvergence,
centering and covariance selection. Source-specific cases cover H2O's
baseline branches, BioJava's restoration, FRESCO's tie disclosure and numeric
range, and JLinAlg's separate diagnostics and backend closure. Compare with
an independently specified finite law or accepted reference implementation,
retaining failed requests in the denominator. Public data then support
conversion and realistic-size trials without pretending to supply known
causal truth. These exact witnesses motivate a prototype and its experiments;
they are not completed prototype evidence. \citep{exec24_longjava_reference}

Full costs include preparation, rights review, environment maintenance,
failed fits, diagnostics, protocol communication and operator time. Wider
search leads, uninspected delegates and native qualification remain explicit.
The comparison replaces an uninformative zero-result search with concrete
evidence and actions, without declaring the area closed.
